% GENERATED FILE. Edit canonical lessons, then run npm run generate:books.
% canonical-source-hash: fnv1a64:510164fcb13720f7
% canonical-lessons: HI-C227-phatna, HI-C227-pharna, HI-C227-palna, HI-C227-cirna, HI-C227-bona

\chapter{To tear, To tear, To bring up, To split, To sow}
\label{ch:hi-to-tear-to-tear-to-bring-up-to-split-to-sow}
\hlchaptermodality{227}

\begin{chapteropening}
\textbf{By the end of this chapter:} \emph{I can say to tear, to tear, to bring up, to split and to sow.}

Complete the last lesson of chapter 227: I can say to tear, to tear, to bring up, to split and to sow.
\end{chapteropening}

\section[\texorpdfstring{\hi{फटना} (phaṭnā)}{फटना (phaṭnā)}]{\hi{फटना} (phaṭnā) — to tear, to burst}

\label{lesson:HI-C227-phatna}



\begin{quote}
Before the new one: say the Hindi for to be made, then the Hindi for to break.
\end{quote}

\subsection*{You'll want to know: \hi{फटना}}
\textbf{\hi{फटना}} — \emph{phaṭnā} — \textquotedblleft{}to tear, to burst\textquotedblright{}.

\begin{cousinweb}
\textbf{\hi{फटना}} (\emph{phaṭnā}) is a tadbhava, a word inherited from Sanskrit and reshaped by sound change. It comes from \emph{sphāṭayati}, \textquotedblleft{}causes to split\textquotedblright{} (CDIAL 13825). \emph{phaṭnā} is the self-acting form, with the short vowel such verbs take.
\end{cousinweb}

\begin{grammarlens}[title={What you've built}]
One more word for this chapter.
\end{grammarlens}

\subsection*{Your turn}
\begin{itemize}
\raggedright
  \item \emph{Say it:} \emph{phaṭnā}
  \item \emph{Say it:} \emph{phaṭnā}, once more
  \item \emph{Say it:} say \emph{phaṭnā}
  \item \emph{Recall:} say the Hindi for to be made, then the Hindi for to break, then say \emph{phaṭnā} again
\end{itemize}

\begin{tcolorbox}[breakable,colback=teal!4,colframe=teal!35!black,title={Before you move on}]
What does \hi{फटना} mean? (\textquotedblleft{}To tear, to burst\textquotedblright{}.) Say it once more.
\end{tcolorbox}
\section[\texorpdfstring{\hi{फाड़ना} (phāṛnā)}{फाड़ना (phāṛnā)}]{\hi{फाड़ना} (phāṛnā) — to tear (something)}

\label{lesson:HI-C227-pharna}



\begin{quote}
Before the new one: say the Hindi for to break, then the Hindi for to tear.
\end{quote}

\subsection*{You'll want to know: \hi{फाड़ना}}
\textbf{\hi{फाड़ना}} — \emph{phāṛnā} — \textquotedblleft{}to tear (something)\textquotedblright{}.

\begin{cousinweb}
\textbf{\hi{फाड़ना}} (\emph{phāṛnā}) is a tadbhava, a word inherited from Sanskrit and reshaped by sound change. It comes from \emph{sphāṭayati}, \textquotedblleft{}causes to split\textquotedblright{} (CDIAL 13825).
\end{cousinweb}

\begin{grammarlens}[title={What you've built}]
One more word for this chapter.
\end{grammarlens}

\subsection*{Your turn}
\begin{itemize}
\raggedright
  \item \emph{Say it:} \emph{phāṛnā}
  \item \emph{Say it:} \emph{phāṛnā}, once more
  \item \emph{Say it:} say \emph{phāṛnā}
  \item \emph{Recall:} say the Hindi for to break, then the Hindi for to tear, then say \emph{phāṛnā} again
\end{itemize}

\begin{tcolorbox}[breakable,colback=teal!4,colframe=teal!35!black,title={Before you move on}]
What does \hi{फाड़ना} mean? (\textquotedblleft{}To tear (something)\textquotedblright{}.) Say it once more.
\end{tcolorbox}
\section[\texorpdfstring{\hi{पालना} (pālnā)}{पालना (pālnā)}]{\hi{पालना} (pālnā) — to bring up, to keep (animals)}

\label{lesson:HI-C227-palna}



\begin{quote}
Before the new one: say the Hindi for to tear, then the Hindi for to tear.
\end{quote}

\subsection*{You'll want to know: \hi{पालना}}
\textbf{\hi{पालना}} — \emph{pālnā} — \textquotedblleft{}to bring up, to keep (animals)\textquotedblright{}.

\begin{cousinweb}
\textbf{\hi{पालना}} (\emph{pālnā}) is a tadbhava, a word inherited from Sanskrit and reshaped by sound change. It comes from \emph{pālayati}, \textquotedblleft{}guards, protects\textquotedblright{} (CDIAL 8129).
\end{cousinweb}

\begin{grammarlens}[title={What you've built}]
One more word for this chapter.
\end{grammarlens}

\subsection*{Your turn}
\begin{itemize}
\raggedright
  \item \emph{Say it:} \emph{pālnā}
  \item \emph{Say it:} \emph{pālnā}, once more
  \item \emph{Say it:} say \emph{pālnā}
  \item \emph{Recall:} say the Hindi for to tear, then the Hindi for to tear, then say \emph{pālnā} again
\end{itemize}

\begin{tcolorbox}[breakable,colback=teal!4,colframe=teal!35!black,title={Before you move on}]
What does \hi{पालना} mean? (\textquotedblleft{}To bring up, to keep (animals)\textquotedblright{}.) Say it once more.
\end{tcolorbox}
\section[\texorpdfstring{\hi{चीरना} (cīrnā)}{चीरना (cīrnā)}]{\hi{चीरना} (cīrnā) — to split, to cut open}

\label{lesson:HI-C227-cirna}



\begin{quote}
Before the new one: say the Hindi for to tear, then the Hindi for to bring up.
\end{quote}

\subsection*{You'll want to know: \hi{चीरना}}
\textbf{\hi{चीरना}} — \emph{cīrnā} — \textquotedblleft{}to split, to cut open\textquotedblright{}.

\begin{cousinweb}
\textbf{\hi{चीरना}} (\emph{cīrnā}) is a tadbhava, a word inherited and reshaped by sound change. It comes from \emph{*cīrayati}, \textquotedblleft{}splits\textquotedblright{}, a form CDIAL reconstructs from its descendants (CDIAL 4844). It is built on \emph{cīra}, \textquotedblleft{}strip of cloth\textquotedblright{}.
\end{cousinweb}

\begin{grammarlens}[title={What you've built}]
One more word for this chapter.
\end{grammarlens}

\subsection*{Your turn}
\begin{itemize}
\raggedright
  \item \emph{Say it:} \emph{cīrnā}
  \item \emph{Say it:} \emph{cīrnā}, once more
  \item \emph{Say it:} say \emph{cīrnā}
  \item \emph{Recall:} say the Hindi for to tear, then the Hindi for to bring up, then say \emph{cīrnā} again
\end{itemize}

\begin{tcolorbox}[breakable,colback=teal!4,colframe=teal!35!black,title={Before you move on}]
What does \hi{चीरना} mean? (\textquotedblleft{}To split, to cut open\textquotedblright{}.) Say it once more.
\end{tcolorbox}
\section[\texorpdfstring{\hi{बोना} (bonā)}{बोना (bonā)}]{\hi{बोना} (bonā) — to sow}

\label{lesson:HI-C227-bona}



\begin{quote}
Before the new one: say the Hindi for to bring up, then the Hindi for to split.
\end{quote}

\subsection*{You'll want to know: \hi{बोना}}
\textbf{\hi{बोना}} — \emph{bonā} — \textquotedblleft{}to sow\textquotedblright{}.

\begin{cousinweb}
\textbf{\hi{बोना}} (\emph{bonā}) goes back to Sanskrit \emph{vapati}, \textquotedblleft{}scatters, sows\textquotedblright{} (CDIAL 11282). CDIAL marks the Hindi form as taken over from Gujarati.
\end{cousinweb}

\begin{grammarlens}[title={What you've built}]
One more word for this chapter.
\end{grammarlens}

\subsection*{Your turn}
\begin{itemize}
\raggedright
  \item \emph{Say it:} \emph{bonā}
  \item \emph{Say it:} \emph{bonā}, once more
  \item \emph{Say it:} say \emph{bonā}
  \item \emph{Recall:} say the Hindi for to bring up, then the Hindi for to split, then say \emph{bonā} again
\end{itemize}

\begin{tcolorbox}[breakable,colback=teal!4,colframe=teal!35!black,title={Before you move on}]
What does \hi{बोना} mean? (\textquotedblleft{}To sow\textquotedblright{}.) Say it once more.
\end{tcolorbox}
